% Generated draft: requires mathematical and source review.
\documentclass[11pt]{article}
\usepackage[margin=25mm]{geometry}
\usepackage{fontspec}
\setmainfont{Latin Modern Roman}
\usepackage{amsmath,amssymb,mathtools,unicode-math}
\setmathfont{Latin Modern Math}
\usepackage{enumitem,needspace,xurl,hyperref}
\setlength{\parindent}{0pt}
\setlength{\parskip}{5pt}
\setlength{\emergencystretch}{3em}
\newcommand{\N}{\mathbb{N}}
\newcommand{\Z}{\mathbb{Z}}
\newcommand{\Q}{\mathbb{Q}}
\newcommand{\R}{\mathbb{R}}
\newcommand{\C}{\mathbb{C}}
\newcommand{\F}{\mathbb{F}}
\newcommand{\A}{\mathbb{A}}
\newcommand{\PP}{\mathbb{P}}
\newcommand{\E}{\mathbb{E}}
\newcommand{\eps}{\varepsilon}
\newcommand{\dd}{\,\mathrm{d}}
\begin{document}
\section*{UnsolvedMath ID: 30002001}
\textbf{Source problem code:} catalogue entry
\subsection{Definitions and mathematical statement}
An affine Weyl group is a Coxeter group arising from an affine root system; write \((W,S)\) for such a group with its simple reflections, and write \(\leq\) for its Bruhat order. Under the usual equal-parameter interpretation, its Hecke algebra \(\mathcal H(W)\) is the algebra over \(\Z[v,v^{-1}]\) with standard basis \(\{T_x:x\in W\}\), the braid relations, and quadratic relations
\[
(T_s+v)(T_s-v^{-1})=0\qquad(s\in S).
\]
The bar involution sends \(v\) to \(v^{-1}\) and \(T_x\) to \(T_{x^{-1}}^{-1}\). In this normalization, the Kazhdan--Lusztig basis element, denoted by \(H_x\) in the supplied record, is characterized by
\[
\overline{H_x}=H_x,
\qquad
H_x\in T_x+\sum_{y<x}v\Z[v]T_y.
\]
For a prime \(p\), the \(p\)-canonical basis is the basis obtained from normalized characters of indecomposable parity complexes in characteristic \(p\) in a compatible geometric realization of the Hecke algebra. Denote its element indexed by \(x\) by \({}^{p}H_x\). Any root datum or realization data needed for this construction must be fixed; denote these collectively by \(D\).

\textbf{Classification problem.} For fixed \((W,S)\) and compatible data \(D\), determine the set
\[
\mathcal C(W,D)=\bigl\{(p,x):p\text{ is prime},\ x\in W,\ {}^{p}H_x=H_x\text{ in }\mathcal H(W)\bigr\}.
\]
Thus the requested criterion should characterize exactly when the equality holds, for every prime \(p\) and every element \(x\) of the chosen affine Weyl group.

\textbf{Unresolved scope.} The supplied record does not specify an affine type, root datum, geometric realization, or restrictions on the characteristic. It also does not explicitly state the Hecke-algebra normalization. The definitions above give a conventional interpretation with compatible normalization of both bases; they do not identify a uniquely specified realization or settle whether the intended classification concerns one fixed affine Weyl group or all affine Weyl groups.
\subsection{Short English statement}
Determine exactly which primes \(p\) and affine Weyl-group elements \(x\) have the same \(p\)-canonical and Kazhdan--Lusztig basis element. The affine Weyl group and any required realization data must be specified before this becomes a fully determined classification problem.
\subsection{Sources}
\begin{itemize}
\item Supplied local UnsolvedMath registry record, permanent ID \(30002001\), source problem code \texttt{catalogue entry}, titled ``Equality of \(p\)-Canonical and Kazhdan--Lusztig Bases,'' in the collection \emph{Oberwolfach Reports Open Problems}. The record assigns the problem to 2012 and names no individual proposer.
\item \emph{Enveloping Algebras and Geometric Representation Theory} (2012), Oberwolfach Reports, pp.\ 733--809, DOI \texttt{10.4171/owr/2012/13}. Supplied source report and literature reference: \url{https://doi.org/10.4171/owr/2012/13}. The supplied background describes these reports as documenting mathematical workshops and their open-problem discussions; no exact problem page or speaker attribution is supplied.
\item \emph{Koszul Duality for Kac-Moody Groups and Characters of Tilting Modules}. Supplied literature reference: \url{https://arxiv.org/abs/1706.00183}. No specific theorem or passage supporting the classification problem is identified in the record.
\item Literature assessment dated \texttt{2026-08-21}: the supplied curation labels the problem open and states that equality is known in substantial ranges and under prime-size hypotheses, while no general classification of all pairs \((p,x)\) has been proved. This is supplied dated curation, not an independently verified literature assessment. The record expressly requests independent verification of time-sensitive status and cited claims.
\item The registry labels the statement \texttt{corrected\_verified} and reports a manual semantic, typographic, and LaTeX review of its cleanup. These are supplied editorial assertions, not independent mathematical or source verification of this draft.
\end{itemize}
\end{document}
